%% ma: L11-B15-0051
%% lop: 11
%% chuong: 5
%% bai: 15
%% dang: 11.15.5
%% loai: TN4
%% muc: VD
%% dap_an: "D"
%% nguon: "tai-lieu-hoc-tap-mon-toan-lop-11-theo-chuong-trinh-giao-duc-pho-thong-2018.pdf (Lớp Toán Thầy Duy, Phạm Lê Duy, THPT Chu Văn An), Chương 3 Bài 1, Phần 1, Câu 40 (tr.329)"
%% kien_thuc: "Lớp 11 Bài 15 (giới hạn của dãy số)"
%% trang_thai: cach-ly
%% ly_do: "Đề và đáp án gốc không khớp: theo đề, lim = 1/2, không có trong bốn phương án; đáp án gốc D (1/6) ứng với căn bậc ba là ∛(n^3+n^2). Gợi ý sửa: đổi n^3+n thành n^3+n^2 (khi đó L = 1/2 − 1/3 = 1/6), hoặc đổi một phương án thành 1/2."
%% ghi_chu: "Câu có lỗi đề/đáp án của tài liệu gốc."
\begin{ex}
Tính giới hạn $L=\lim\left(\sqrt{n^2+n+1}-\sqrt[3]{n^3+n}\right)$.
\choice{$+\infty$}{$\dfrac54$}{$\dfrac{53}2$}{\True $\dfrac16$}
\loigiai{\sloppy $\sqrt{n^2+n+1}-n=\dfrac{n+1}{\sqrt{n^2+n+1}+n}\to\dfrac12$ và $\sqrt[3]{n^3+n}-n=\dfrac{n}{B^2+nB+n^2}\to0$ với $B=\sqrt[3]{n^3+n}$. Vậy $L=\dfrac12$, không có trong các phương án.}
\end{ex}
